%------------------------------------------------------------------------------%
\begin{question}
  Determine as soluções do sistema
  \[
    \systeme{
       x + 2y = 0,
      3x -  y = 0
    }
  \]

  \bigskip
  \pause
  Matriz aumentada
  \[
    \begin{amatrix}{2}
      1 & 2  & 0 \\
      3 & -1 & 0
    \end{amatrix}
  \]
\end{question}

%------------------------------------------------------------------------------%
\begin{resolution}{Escalonamento}
  \begin{columns}
  \column{0.3\linewidth}
    \[\;\;\;
    \left[\begin{array}{cc|c}
      1 & 2  & 0 \\
      3 & -1 & 0
    \end{array}\right]
    \]
    \pause
  \column{0.7\linewidth}
    Eliminamos os elementos abaixo do primeiro pivô
    \\ \pause
    \(
      L_2 \leftarrow L_2 - 3 L_1
    \)
  \end{columns}
  \vspace{-\baselineskip}
  \[
    \pause
    L'_2
    \;=\;
    L_2 - 3 L_1
    \pause
    \;=\;
    \left[\begin{array}{cc|c}
      3 & -1 & 0
    \end{array}\right]
    - 3
    \left[\begin{array}{cc|c}
      1 & 2  & 0
    \end{array}\right]
    \pause
    \;=\;
    \left[\begin{array}{cc|c}
      0 & -7 & 0
    \end{array}\right]
  \]
  \pause
  \[
    \left[\begin{array}{cc|c}
      1 & 2  & 0 \\
      0 & -7 & 0
    \end{array}\right]
  \]
\end{resolution}

%------------------------------------------------------------------------------%
\begin{resolution}{Escalonamento}
  \begin{columns}
  \column{0.3\linewidth}
    \[\;\;\;
    \left[\begin{array}{cc|c}
      1 & 2  & 0 \\
      0 & -7 & 0
    \end{array}\right]
    \]
    \pause
  \column{0.7\linewidth}
    Tornando o segundo pivô igual a 1
    \\ \pause
    \(
      L_2 \leftarrow - \frac{1}{7} L_2
    \)
  \end{columns}
  \vspace{-\baselineskip}
  \[
    \pause
    L'_2
    \;=\;
    - \frac{1}{7} L_2
    \pause
    \;=\;
    - \frac{1}{7}
    \left[\begin{array}{cc|c}
      0 & -7 & 0
    \end{array}\right]
    \pause
    \;=\;
    \left[\begin{array}{cc|c}
      0 & 1 & 0
    \end{array}\right]
  \]
  \pause
  \[
    \left[\begin{array}{cc|c}
      1 & 2 & 0 \\
      0 & 1 & 0
    \end{array}\right]
  \]
\end{resolution}

%------------------------------------------------------------------------------%
\begin{resolution}{Escalonamento}
  \begin{columns}
  \column{0.3\linewidth}
    \[\;\;\;
    \left[\begin{array}{cc|c}
      1 & 2 & 0 \\
      0 & 1 & 0
    \end{array}\right]
    \]
    \pause
  \column{0.7\linewidth}
    Eliminamos os elementos a cima do segundo pivô
    \\ \pause
    \(
      L_1 \leftarrow L_1 - 2 L_2
    \)
  \end{columns}
  \vspace{-\baselineskip}
  \[
    \pause
    L'_1
    \;=\;
    L_1 - 2 L_2
    \pause
    \;=\;
    \left[\begin{array}{cc|c}
      1 & 2 & 0
    \end{array}\right]
    - 2
    \left[\begin{array}{cc|c}
      0 & 1 & 0
    \end{array}\right]
    \pause
    \;=\;
    \left[\begin{array}{cc|c}
      1 & 0 & 0
    \end{array}\right]
  \]
  \pause
  \[
    \left[\begin{array}{cc|c}
      1 & 0 & 0 \\
      0 & 1 & 0
    \end{array}\right]
  \]

  \pause
  \bigskip
  Observe que o a última coluna da matriz aumentada sempre ficou igual a zero
\end{resolution}

%------------------------------------------------------------------------------%
\begin{resolution}{Solução}
  A única solução é
  \begin{align*}
    x & = 0 \\
    y & = 0
  \end{align*}

  \pause
  Vetorialmente
  \[
    \mathbf{x}
    \pause
    =
    \begin{bmatrix}
      x \\
      y
    \end{bmatrix}
    \pause
    =
    \begin{bmatrix}
      0 \\
      0
    \end{bmatrix}
  \]
\end{resolution}

%------------------------------------------------------------------------------%
\endinput
